\chapter{Konzeption des Formular-Standards und der Prüf-Registry}
\label{chap:konzeption}

Aufbauend auf den Befunden der empirischen Ist-Analyse wird in diesem Kapitel der formale Formular-Standard für die Hochschule für Wirtschaft und Gesellschaft Ludwigshafen entwickelt. Dieser Standard adressiert nicht nur ergonomische und gestalterische Dimensionen nach DIN EN ISO 9241-110 und BITV 2.0, sondern integriert auch rechtliche Vorgaben der Datenschutz-Grundverordnung (DSGVO). Anschließend wird dieser Standard in eine maschinenlesbare Komponenten- und Regel-Registry überführt, die als deterministische Wissensbasis für die automatisierte Compliance-Prüfung dient.

\section{Struktur-, Navigations- und Informationsarchitektur}
\label{sec:architektur}

Eine verlässliche Formular-Governance bedarf klar definierter architektonischer Leitplanken. Die heterogene Verteilung zwischen 700px und 850px im Bestand wird durch eine verbindliche Standardbreite harmonisiert:

\begin{itemize}
    \item \textbf{Responsives Desktop-Raster (850px):} Alle Formularvorlagen werden verbindlich auf eine Maximalbreite von \texttt{850px} (\texttt{max-width: 850px; min-width: 300px;}) normiert. Dies garantiert einen ruhigen Zeilenrhythmus, verhindert übermäßige Zeilenlängen bei Textfeldern und gewährleistet eine konsistente Einbindung in das Hochschulportal.
    \item \textbf{Seitenstrukturierung und Pagination (Chunking):} Formulare mit mehr als 20 Einzelelementen müssen in thematisch abgegrenzte Seiten oder Abschnitte unterteilt werden. Für mehrseitige Formulare wird ein linearer Pagination-Flow vorgegeben.
    \item \textbf{Navigationsschaltflächen:} Schaltflächen zur Seitensteuerung müssen strikt mit den handlungsorientierten Verben \textit{\enquote{Weiter}} und \textit{\enquote{Zurück}} bezeichnet werden. Technische Bezeichnungen wie \textit{\enquote{weiter zu Seite 2}} sind unzulässig.
\end{itemize}

\section{Corporate Design, Barrierefreiheit und Stammdaten-Muster}
\label{sec:cd_komponenten}

Die visuelle und programmatische Gestaltung bildet das Bindeglied zwischen institutioneller Wiedererkennung und digitaler Barrierefreiheit.

\subsection{Corporate Identity, Header- und Footer-Struktur}
Jedes Formular muss im Kopfbereich über das offizielle Bildbanner der Hochschule (\texttt{XImage}) verfügen. Zur Wahrung der Barrierefreiheit (WCAG 2.1 Erfolgskriterium 1.1.1) ist ein textueller Alternativtext zwingend vorgeschrieben (\enquote{Hochschule für Wirtschaft und Gesellschaft Ludwigshafen}).

Darüber hinaus muss im Fußbereich (\texttt{XFooter}) der herstellerseitige Werbehinweis \textit{\enquote{Powered by XIMA Formcycle}} deaktiviert werden, indem im Formcycle-Designer die Eigenschaft \textit{\enquote{Logo anzeigen}} (\texttt{properties.showlogo = "0"}) abgewählt wird (\texttt{HWG-CD-03}). Dies verhindert fremdes Software-Branding und stärkt die institutionelle Seriosität.

\subsection{Semantische Gruppierung und Pflichtfeld-Indikation}
Gemäß WCAG 2.1 Erfolgskriterium 1.3.1 (Info und Beziehungen) dürfen Container nicht ausschließlich über visuelle \texttt{<div>}-Tags realisiert werden. Jeder logische Block (z.\,B. persönliche Angaben, Anschrift, Bankdaten) muss als semantisches \texttt{<fieldset>} mit einer lesbaren Legende (\texttt{<legend>}) ausgezeichnet werden (\texttt{HWG-A11Y-01}). 

Zur Wahrung des Erfolgskriteriums 3.3.2 (Beschriftungen) muss jedes Pflichtfeld sowohl optisch über einen konsistenten Pflichtfeld-Stern (\texttt{showrequiredhint="1"}) als auch maschinenlesbar im DOM über \texttt{aria-required="true"} deklariert sein (\texttt{HWG-A11Y-02}). Technisch erzeugt Formcycle das HTML-Attribut \texttt{aria-required="true"} automatisch im gerenderten DOM, sobald im Formcycle-Designer unter \textit{Einschränkungen} das Kontrollkästchen \textit{\enquote{Pflichtfeld}} (\texttt{required="1"}) gemeinsam mit \textit{\enquote{Pflichtfeldhinweis anzeigen}} (\texttt{showrequiredhint="1"}) aktiviert wird. Wird ein Feld hingegen lediglich dynamisch per Bedingung (\texttt{requiredif}) als erforderlich geschaltet, muss \texttt{aria-required="true"} explizit im Designer-Reiter \textit{Erweitert $\rightarrow$ HTML-Attribute} (\texttt{properties.attributes}) hinterlegt werden.

\subsection{Das normierte HWG-Stammdaten-Muster}
Um die im Bestand identifizierten Copy-Paste-Abweichungen zu eliminieren, wird die Erfassung persönlicher Angaben standardisiert:
\begin{enumerate}
    \item \textbf{Feste Sequenz (\texttt{HWG-STAMM-01}):} Die Erfassung folgt dem natürlichen Lesefluss: $\text{Anrede} \rightarrow \text{Vorname} \rightarrow \text{Nachname} \rightarrow \text{E-Mail-Adresse}$. Eine Voranstellung des Nachnamens ist nur in behördlichen Spezialregistern gestattet.
    \item \textbf{Kanonische Stammdaten-Bezeichnungen (\texttt{HWG-STAMM-02}):} Sämtliche Stammdatenfelder folgen einer einheitlichen Terminologie (\textit{\enquote{Anrede}}, \textit{\enquote{Vorname}}, \textit{\enquote{Nachname}}, \textit{\enquote{E-Mail-Adresse}} bzw. \textit{\enquote{Hochschul-E-Mail}}, \textit{\enquote{Matrikelnummer}}, \textit{\enquote{Postleitzahl}}, \textit{\enquote{Wohnort}}, \textit{\enquote{IBAN}}). Redundante Rollenzusätze im Label (wie \enquote{Name Antragstellende/r}) oder unorthodoxe Schreibweisen (\enquote{HWGLU-Email}) sind unzulässig.
    \item \textbf{Bereinigung von Entwickler-Rückständen (\texttt{HWG-UI-01}):} Sichtbare technische Implementierungs-Notizen (z.\,B. Verweise auf JavaScript, Shibboleth-Gruppen oder LDAP) sind im Produktivzustand strikt untersagt.
\end{enumerate}

\section{Datenschutz, Eingabevalidierung und Schaltflächen-Semantik}
\label{sec:datenschutz_validierung}

Über das Layout hinaus werden zentrale funktionale Anforderungen definiert:

\subsection{Datenschutz-Transparenz (Art. 13 DSGVO)}
Sobald ein Formular personenbezogene Daten verarbeitet, muss vor der Absende-Schaltfläche oder im aktivierten Datenschutz-Bereich des Footers (\texttt{privacyPolicyShow="1"}) ein verständlicher Hinweis auf den Verarbeitungszweck sowie ein klickbarer Hyperlink auf die zentrale Datenschutzerklärung der HWG Ludwigshafen hinterlegt sein (\texttt{HWG-DSGVO-01}).

\subsection{Strukturierte Validierung, Dateiuploads und Zwischenspeichern-Empfehlung}
Eingabefelder für strukturierte Datenformate (IBAN, Postleitzahl, Telefonnummer, Datum) dürfen nicht als unvalidierte Freitexte verbleiben (\texttt{HWG-VAL-01}). Es sind formale Validierungsmuster (Regex) oder Formcycle-Presets zu hinterlegen. Kurze Felder wie die fünfstellige Postleitzahl oder Hausnummern dürfen nicht auf 100\,\% Zeilenbreite gestreckt werden, um dem Nutzenden die erwartete Inhaltslänge visuell zu signalisieren (\texttt{HWG-UX-01}, Affordanz nach Nielsen Norman Group).

Für Dateiuploads (\texttt{XUpload}, \texttt{XFile}) gilt (\texttt{HWG-UPLOAD-01}): Das Label oder der Hilfetext muss die akzeptierten Dateiformate (z.\,B. \enquote{Nur PDF, JPG}) sowie die maximale Dateigröße (z.\,B. \enquote{max. 10 MB}) transparent ausweisen.

Überschreitet ein Formular eine kritische Komplexitätsschwelle ($\ge 3$ Formularseiten \texttt{XPage} oder $\ge 25$ Eingabefelder), spricht der Standard eine unverbindliche Gestaltungsempfehlung (\texttt{HWG-UX-02}) ohne Punktabzug aus: Es sollte eine \textit{\enquote{Zwischenspeichern}}-Schaltfläche angeboten werden, die den Bearbeitungsstand sichert und den Nutzenden automatisch einen personalisierten Fortsetzungs-Link per E-Mail zusendet.

\subsection{Schaltflächen-Semantik und Workflow-E-Mail-Transparenz}
Aktionsschaltflächen müssen zwingend dem \textit{Action-Verb-Muster} folgen, also eine Kombination aus Zielobjekt und Handlungsverb im Infinitiv aufweisen (z.\,B. \textit{\enquote{Antrag einreichen}} oder \textit{\enquote{Anfrage absenden}} statt dem kleingeschriebenen generischen \textit{\enquote{senden}}, \texttt{HWG-BTN-01}). Zustandsbezeichnungen und fachliche Ausnahme-Workflows (wie die Fallstudie \textit{\enquote{Kritischer Fall}}) dürfen niemals als bloße Substantive auf einer Schaltfläche stehen, sondern müssen die ausgelöste Aktion explizit benennen (z.\,B. \textit{\enquote{Kritischen Fall an Präsidenten weiterleiten}}).

Zur Sicherstellung der Prozesstransparenz nach dem Absenden (\texttt{HWG-FEEDBACK-01}) wird nicht das statische Formular-Frontend, sondern die im Formcycle-Projektarchiv (\texttt{.fcexport}) hinterlegte Workflow-Definition geprüft: In den E-Mail-Knoten (\texttt{type="FC\_EMAIL"}) muss eine automatisierte Eingangsbestätigung an die antragstellende Person konfiguriert sein, die den Eingang bestätigt und über die weiteren Prüfschritte informiert.

\section{Formalisierung als maschinenlesbare Komponenten- und Regel-Registry}
\label{sec:registry_definition}

Um die Einhaltung dieser Vorgaben automatisiert überprüfen zu können, werden alle Kriterien in ein maschinenlesbares JSON-Schema (\texttt{rules\_registry.json}, Version 1.3.0) überführt.

\subsection{Mathematisches Bewertungsmodell}
Das Konformitätsniveau eines Formulars wird über einen normalisierten Compliance-Score $S \in [0, 100]$ quantifiziert. Ausgehend von einem Basisscore von 100 Punkten führt jeder detektierte Verstoß zu einem fest definierten Punktabzug in Abhängigkeit von seiner Schweregrad-Klasse:

\begin{equation}
    S = \max\left(0, \, 100 - \sum_{i=1}^{M} w(r_i) \cdot v(r_i)\right)
    \label{eq:compliance_score}
\end{equation}

Hierbei bezeichnet:
\begin{itemize}
    \item $v(r_i) \in \{0, 1\}$ den binären Auslösestatus der Regel $r_i$ ($1 = \text{Verstoß bzw. Hinweis ausgelöst}$, $0 = \text{konform}$).
    \item $w(r_i)$ das gewichtete Strafmaß der jeweiligen Schweregrad-Kategorie:
    \begin{equation}
        w(r_i) = \begin{cases}
            15 & \text{für } \textit{critical} \text{ (rechtliche/technische Blocker)} \\
            8  & \text{für } \textit{major} \text{ (schwere Usability-/CI-Verletzungen)} \\
            3  & \text{für } \textit{minor} \text{ (visuelle oder syntaktische Mängel)} \\
            0  & \text{für } \textit{recommendation} \text{ (unverbindliche UX-Empfehlung ohne Abzug)}
        \end{cases}
    \end{equation}
\end{itemize}

\subsection{Katalog der formalisierten Governance-Regeln}
Tabelle~\ref{tab:rules_overview} fasst die 19 formalisierten Prüfregeln (18 bewertete Compliance-Regeln sowie 1 unverbindliche Empfehlungsregel) der Registry v1.3.0 zusammen:

\begin{table}[H]
\centering
\scriptsize
\begin{tabularx}{\textwidth}{llllX}
\toprule
\textbf{Regel-ID} & \textbf{Bezeichnung} & \textbf{Kategorie} & \textbf{Typ} & \textbf{Schweregrad \& Begründung} \\
\midrule
\texttt{HWG-CD-01} & Standardbreite 850px & Corporate Design & Deterministisch & Major (8 Pkt): Einheitliches Raster \\
\texttt{HWG-CD-02} & Offizieller CI-Header & Corporate Design & Deterministisch & Critical (15 Pkt): Authentizität / Logo \\
\texttt{HWG-CD-03} & Kein XIMA-Logo im Footer & Corporate Design & Deterministisch & Major (8 Pkt): \texttt{showlogo="0"} im Footer \\
\texttt{HWG-A11Y-01} & Fieldset-Gruppierung & Accessibility & Deterministisch & Critical (15 Pkt): WCAG 1.3.1 \\
\texttt{HWG-A11Y-02} & Pflichtfeld-Stern/ARIA & Accessibility & Deterministisch & Major (8 Pkt): WCAG 3.3.2 (\texttt{required="1"}) \\
\texttt{HWG-A11Y-03} & Saubere Label-Syntax & Accessibility & Deterministisch & Minor (3 Pkt): Keine HTML-Entities \\
\texttt{HWG-STAMM-01} & Stammdaten-Sequenz & Stammdaten & Deterministisch & Major (8 Pkt): Natürlicher Lesefluss \\
\texttt{HWG-STAMM-02} & Kanonische Stammdaten & Stammdaten & Semantisch & Minor (3 Pkt): Einheitliche Feldlabels \\
\texttt{HWG-UI-01} & Keine Dev-Notizen & UI-Qualität & Deterministisch & Critical (15 Pkt): Saubere UI ohne Artefakte \\
\texttt{HWG-BTN-01} & Action-Verb-Muster Buttons & Microcopy & Semantisch & Major (8 Pkt): Zielobjekt + Vollverb \\
\texttt{HWG-BTN-02} & Einheitliche Pagination & Microcopy & Semantisch & Minor (3 Pkt): Weiter/Zurück \\
\texttt{HWG-DSGVO-01} & Datenschutz-Transparenz & Datenschutz & Deterministisch & Critical (15 Pkt): Art. 13 DSGVO \\
\texttt{HWG-UPLOAD-01} & Upload-Spezifikation & Validierung/UX & Deterministisch & Major (8 Pkt): Dateitypen \& Max-Größe \\
\texttt{HWG-VAL-01} & Formatvalidierung & Validierung/UX & Deterministisch & Major (8 Pkt): IBAN/PLZ/Datum-Regex \\
\texttt{HWG-UX-01} & Feldbreiten-Affordanz & Validierung/UX & Deterministisch & Minor (3 Pkt): Kompakte Breiten \\
\texttt{HWG-UX-02} & Zwischenspeichern-Option & Validierung/UX & Semantisch & Recommendation (0 Pkt): Ab 3 S. / 25 Feldern \\
\texttt{HWG-LANG-01} & Amtsjargon mit Erklärung & Leichte Sprache & Semantisch & Major (8 Pkt): WCAG 3.1 \\
\texttt{HWG-LANG-02} & Barrierearme Sprache & Leichte Sprache & Semantisch & Minor (3 Pkt): Keine Schachtelsätze \\
\texttt{HWG-FEEDBACK-01} & Workflow-E-Mail-Feedback & Transparenz & Semantisch & Minor (3 Pkt): \texttt{FC\_EMAIL} im Workflow \\
\bottomrule
\end{tabularx}
\caption{Übersicht der formalisierten Formular-Governance-Regeln in der Registry v1.3.0}
\label{tab:rules_overview}
\end{table}

Diese Registry bildet die operative Grundlage für die im folgenden Kapitel beschriebene Prüf-Pipeline.
